% Evaluation section, v3. Included from paper.tex with % Evaluation section, v3. Included from paper.tex with % Evaluation section, v3. Included from paper.tex with % Evaluation section, v3. Included from paper.tex with \input{evaluation_v3}.
%
% Bibliography keys required:
%   pangakis2023validation, krippendorff2013content, krippendorff2016unitizing,
%   mathet2015gamma, mathet2017gammacat, luo2005coreference,
%   kirillov2019panoptic, chinchor1992muc4, chen2023unified, dieugenio2004kappa
%
% Figures are computed after validating every label in both files against the
% codebook vocabulary (one variant spelling in the gold file, "asymetrical",
% was silently costing seven points on reciprocity before this check existed).

\section{Evaluating the annotation pipeline}
\label{sec:evaluation}

\subsection{A gold standard of one novel}

Before running the codebook over some four hundred novels we need to know what
an LLM annotation of a novel is actually worth; automated annotation without
that check is not a measurement \cite{pangakis2023validation}. We therefore
hand-annotated one novel, Lucie Delarue-Mardrus's \emph{L'Ange et les pervers}
(1930), following the same codebook we give the models. The choice is
deliberate rather than convenient: its protagonist, Marion / Mario de
Valdeclare, is presented as intersex and moves between sapphic salons and male
homosexual circles, so the novel exercises the parts of the grid we care most
about --- the \texttt{queer} value of \texttt{relation\_type}, collective and
milieu relationships, and desire that is coded rather than declared. The human
annotation records fourteen relationships, from the central bond between Marion
and Laurette Wells down to a two-line marriage recalled in the backstory. One
novel gives a first reading of how the pipeline behaves, not a corpus-level
estimate of its accuracy; everything below should be read as a pilot.

\begin{table}[h]
  \centering
  \small
  \begin{tabular}{ll}
    \toprule
    Provider & Model \\
    \midrule
    DeepSeek  & V4 Pro \\
    Anthropic & Claude Sonnet 5 \\
    Anthropic & Claude Opus 5 \\
    OpenAI    & GPT-5.6 Luna Pro \\
    OpenAI    & GPT-5.6 Terra \\
    \bottomrule
  \end{tabular}
  \caption{The five models compared, all called through OpenRouter with an
  identical system prompt, codebook, and JSON schema.}
  \label{tab:models}
\end{table}

\subsection{Two measures and their product}

Reading a novel with this grid means doing two rather different things. The
annotator first draws up a \emph{census}: which relationships in this novel are
worth an entry at all? Then, for each entry, they \emph{code} it, choosing one
label per closed field. A model can do one well and the other badly --- find
exactly the fourteen relationships a human found and misread the tone of every
one of them, or code impeccably three relationships nobody would have listed. A
single accuracy figure hides that difference, which is precisely the difference
we want to report.

The distinction is not ours. Content analysis has long separated
\emph{unitizing} --- deciding what the units of analysis are --- from
\emph{categorizing} them, and holds that the reliability of the first cannot be
folded into the reliability of the second
\cite{krippendorff2013content,krippendorff2016unitizing}. The measure that
addresses both at once is $\gamma$, which aligns annotators' units and scores
agreement on that alignment, with $\gamma_{cat}$ isolating the categorization
component \cite{mathet2015gamma,mathet2017gammacat}. We would use $\gamma$ if we
could. We cannot, for two reasons. $\gamma$ assumes units positioned on a
continuum and counts positional discrepancy as part of its dissimilarity,
whereas a relationship is not located at a place in the novel: it is identified
by \emph{who is in it}, and its footprint is the whole book. And our object is
not agreement between interchangeable observers but the distance between a
system and a reference, which is the setting precision and recall were made for.
What follows therefore keeps $\gamma$'s architecture --- align first, score
categorization on the alignment --- on a set-identity criterion rather than a
positional one.

\paragraph{Matching.} Let $G$ be the set of relationships in the human
annotation of a novel and $P$ the set proposed by a model. Neither is an ordered
list, so nothing can be compared until the two are aligned. We align on the
participants: each relationship $r$ carries an identifier $\mathrm{id}(r)$ built
from the characters involved --- \texttt{laurette\_aimee} for the affair between
Laurette Wells and Aimée de Lagres, \texttt{pere\_mere} for Marion's parents ---
recorded in both files, with aliases resolved beforehand (Marion, Mario and Miss
Marion Hervin are one character). Two relationships match when their identifiers
are equal:
%
\begin{align}
M = \{(g,p) \in G \times P \;:\; \mathrm{id}(g) = \mathrm{id}(p)\} .
\end{align}
%
In the vocabulary of coreference scoring this is CEAF under Luo's strictest
similarity $\phi_1$, which credits an alignment only when the two sets are
identical \cite{luo2005coreference}. Luo sets $\phi_1$ aside as too severe for
coreference, where exact agreement on a mention chain is vanishingly rare; it is
usable here only because participant lists are short and canonicalised upstream.
The gain is that the census score carries no similarity function and no
threshold, and a disagreement about who is involved in a relationship is treated
as what it is: a disagreement about the novel, not a spelling problem.

\paragraph{Recognition quality: did the model see the same relationships?} From
$M$ we take precision, recall and their harmonic mean, computed over
relationships rather than tokens:
%
\begin{align}
\mathrm{P} = \frac{|M|}{|P|}, \qquad
\mathrm{R} = \frac{|M|}{|G|}, \qquad
\mathrm{RQ} = \frac{2\,\mathrm{P}\,\mathrm{R}}{\mathrm{P} + \mathrm{R}} .
\end{align}
%
Precision asks whether the relationships a model proposes are ones a human
reader would also have listed; recall asks how much of the novel's relational
fabric it leaves out. RQ says nothing about how well anything was described.

\paragraph{Coding quality: did it fill the grid the same way?} For a matched
pair $(g,p)$ we compare the closed fields one by one and take the proportion on
which the human and the model wrote the same label:
%
\begin{align}
a(g,p) = \frac{1}{K} \sum_{k=1}^{K} \mathbf{1}\!\left[\, g_k = p_k \,\right] ,
\qquad
\mathrm{CQ} = \frac{1}{|M|} \sum_{(g,p) \in M} a(g,p) .
\end{align}
%
Credit is strict: the codebook offers no principled way to call
\texttt{decadent\_sublime} half-right when the annotator wrote
\texttt{moralized\_pathologized}. Every field weighs the same. The grid has
eleven closed fields, but the gold file predates \texttt{explicitness\_in\_plot},
added in schema v2, so this pilot scores $K = 10$. \texttt{confidence} is logged
by the models but is metadata about the annotation, not a property of the
relationship, and is not scored.

\paragraph{Annotation quality: one number, and why it is a product.} Suppose we
scored the whole annotation at once, with a matched relationship counting not as
$1$ but as $a(g,p)$. The harmonic mean of the resulting precision and recall is
exactly $\mathrm{RQ} \times \mathrm{CQ}$: the mean coding score factors out of
the numerator and cancels in the denominator. The product is therefore not a
weighting we chose but that single figure, read in its two parts:
%
\begin{align}
\mathrm{AQ} = \mathrm{RQ} \times \mathrm{CQ} .
\end{align}
%
A model that finds seven tenths of the relationships and gets seven tenths of
the grid right on those has produced about half the annotation we wanted, which
is what $\mathrm{AQ} \approx 0.51$ says. AQ is bounded above by RQ and reaches
it only if every matched relationship is coded perfectly. We always report the
triple $(\mathrm{RQ}, \mathrm{CQ}, \mathrm{AQ})$: two models can arrive at the
same AQ by opposite routes, one prolific and careless, the other cautious and
narrow, and for a corpus study those are not interchangeable.

This is the shape of panoptic quality in image segmentation, where a recognition
quality --- the same $F_1$ over matched objects --- multiplies a segmentation
quality averaging overlap on the matches \cite{kirillov2019panoptic}, and of
template-filling evaluation since MUC \cite{chinchor1992muc4}. We claim no new
measure: RQ, CQ and AQ instantiate a well-described family, match substructures
then normalise \cite{chen2023unified}, for a hermeneutic coding grid. What the
instantiation buys is the ability to read the two capacities apart.

\subsection{What the pipeline gets right, and where it stops}

\begin{table}[h]
  \centering
  \small
  \begin{tabular}{lcccccc}
    \toprule
    Model & $|P|$ & P & R & RQ & CQ & AQ \\
    \midrule
    \textbf{Claude Opus 5} & 14 & 0.71 & 0.71 & 0.71 & 0.72 & \textbf{0.51} \\
    GPT-5.6 Terra          &  6 & 1.00 & 0.43 & 0.60 & 0.78 & 0.47 \\
    Claude Sonnet 5        &  6 & 1.00 & 0.43 & 0.60 & 0.73 & 0.44 \\
    GPT-5.6 Luna Pro       &  5 & 1.00 & 0.36 & 0.53 & 0.66 & 0.35 \\
    DeepSeek V4 Pro        &  7 & 0.86 & 0.43 & 0.57 & 0.47 & 0.27 \\
    \bottomrule
  \end{tabular}
  \caption{Census, coding and combined score for the five models on
  \emph{L'Ange et les pervers}, against fourteen human-annotated relationships,
  sorted by AQ. $|P|$ is the number of relationships the model proposed.}
  \label{tab:model-results}
\end{table}

The $|P|$ column is the most telling. Claude Opus 5 proposes exactly fourteen
relationships and gets ten of them right; the other four propose five to seven
and are almost never wrong about the ones they propose. That second behaviour is
not caution about \emph{coding} but a narrower idea of what counts as a
relationship: those models list the named couples and stop. What they leave out
is the collective material --- Laurette Wells surrounded by her favourites, the
theatre milieu around Ginette Lobre, the powdered young men of Julien Midalge's
circle. For a study of sapphic and homosexual sociability those are not
decorative background; they are the milieux the corpus was assembled to
describe. A model with perfect precision that returns only the foreground
couples would quietly reshape the corpus into a series of couples, which is why
precision alone is a poor guide to model choice here.

\begin{table}[h]
  \centering
  \small
  \begin{tabular}{lcccccc}
    \toprule
    Field & Opus & Terra & Sonnet & Luna & DeepS. & mean \\
    \midrule
    \texttt{relation\_type} & 0.90 & 1.00 & 1.00 & 1.00 & 0.83 & 0.95 \\
    \texttt{structure}      & 0.90 & 1.00 & 0.67 & 1.00 & 1.00 & 0.91 \\
    \texttt{centrality}     & 0.90 & 0.83 & 1.00 & 0.80 & 0.67 & 0.84 \\
    \texttt{status}         & 0.70 & 1.00 & 0.83 & 0.80 & 0.83 & 0.83 \\
    \midrule
    \texttt{explicitness}   & 0.70 & 1.00 & 0.67 & 0.80 & 0.33 & 0.70 \\
    \texttt{reciprocity}    & 0.70 & 0.83 & 0.83 & 0.80 & 0.33 & 0.70 \\
    \texttt{outcome}        & 0.70 & 0.83 & 0.67 & 0.60 & 0.33 & 0.63 \\
    \midrule
    \texttt{exclusivity}    & 0.70 & 0.50 & 0.83 & 0.40 & 0.00 & 0.49 \\
    \texttt{adultery}       & 0.70 & 0.33 & 0.67 & 0.40 & 0.17 & 0.45 \\
    \texttt{framing}        & 0.30 & 0.50 & 0.17 & 0.00 & 0.17 & 0.23 \\
    \bottomrule
  \end{tabular}
  \caption{Agreement with the human annotation, field by field, on matched
  relationships only (10 pairs for Opus, 6 for Terra, Sonnet and DeepSeek, 5 for
  Luna); the last column is the unweighted mean of the five rates. Fields are
  grouped by mean agreement, not by codebook family.}
  \label{tab:field-results}
\end{table}

Averaging over fields hides the finding, so Table~\ref{tab:field-results}
reports agreement field by field, and the grid falls into three layers that are
not the codebook's three families. The top layer is descriptive and essentially
solved: once a model and a human agree that a relationship exists, they agree on
its gender configuration, the shape of the group, the kind of bond and how much
room the novel gives it. The \texttt{queer} coding is the strongest case --- every
relationship involving Marion / Mario is coded \texttt{queer} by all five models,
which is not trivial for a 1930 novel that never uses a clinical vocabulary for
its protagonist.

The middle layer sits at roughly two thirds, usable with care. The bottom layer
is where the pipeline fails, and its three fields have something in common: they
all ask for a judgment about norms. \texttt{exclusivity} and \texttt{adultery}
record the rules a relation operates under and whether they are broken, and the
codebook binds them by an explicit pairing rule; the models do not apply it as a
rule, which suggests a schema constraint rather than a better prompt.
\texttt{framing} resists most. \texttt{neutral\_ambivalent} is the models'
default, given fifteen times out of thirty-three, and
\texttt{redemptive\_spiritualized} --- which the annotator uses five times --- is
never once produced by any model. The vocabulary of \emph{pervers} and of
salvation that gives the novel its title is exactly what the automatic pass does
not recover. That contrast \emph{is} the result: the grid contains a descriptive
layer that can be automated at corpus scale today and a normative layer that
cannot, and the second carries the interpretive claims of the project.

\subsection{Limits of the measure}

Three limits bear directly on how the table should be read. \emph{Exact keying
makes collective relationships unscoreable.} Alias resolution settles naming for
named characters; it does not decide whether Laurette and her favourites are one
\texttt{hierarchy\_group} relationship or six dyads. That is a scholarly
decision, not noise, and under an exact key a legitimate disagreement about it
costs a false positive and a false negative at once. Fixing a canonical
identifier per milieu is the first change planned for the grid, alongside a
\texttt{direction} field, since the participant list stores no roles and a model
that inverts the seducer and the seduced currently scores a perfect match.

\emph{CQ is measured on a sample each model chooses for itself.} Coding quality
is computed only on matched pairs, so Terra's 0.78 averages over six legible
couples while Opus's 0.72 averages over ten, including the collective forms the
others declined --- the prolific model is penalised for attempting the hard
cases. The ordering of the middle three models in
Table~\ref{tab:model-results} should not be read as a separation.

\emph{Agreement is not chance-corrected, and there is no human ceiling.} The
marginals are skewed and the models almost never abstain where the annotator
does, so the top of the field table is flattered \cite{dieugenio2004kappa}. A
chance-corrected coefficient becomes meaningful once a second human annotator
gives us something to compare it against: fields where two readers would agree
at 0.6 cannot be used to convict a model at 0.6, and double annotation of a
subset is the measurement that would make
Table~\ref{tab:field-results} fully interpretable.

\subsection{What we carry forward}

On this pilot, Claude Opus 5 is the model we take to the rest of the corpus: it
is the only one that annotates at the granularity the codebook asks for, milieux
included, and it pays for that coverage with a precision the others never have
to risk. The ranking rests on one novel and should be revisited as the gold
standard grows.

More important than the ranking is the shape of the result. The census and the
descriptive half of the grid --- who is involved, in what gender configuration,
in what shape of group, in what kind of bond, with how much narrative room ---
are reliable enough to run over four hundred novels and to support
distributional claims. The normative half --- what the relation owes to whom,
and what the novel makes of it --- is not, and we do not intend to publish
corpus-wide figures on \texttt{framing} on the strength of an automatic pass.
The pipeline's role is a division of labour rather than a replacement: it
produces the census and the descriptive coding at scale, and directs human
reading towards the relationships where the interpretive fields are actually at
stake..
%
% Bibliography keys required:
%   pangakis2023validation, krippendorff2013content, krippendorff2016unitizing,
%   mathet2015gamma, mathet2017gammacat, luo2005coreference,
%   kirillov2019panoptic, chinchor1992muc4, chen2023unified, dieugenio2004kappa
%
% Figures are computed after validating every label in both files against the
% codebook vocabulary (one variant spelling in the gold file, "asymetrical",
% was silently costing seven points on reciprocity before this check existed).

\section{Evaluating the annotation pipeline}
\label{sec:evaluation}

\subsection{A gold standard of one novel}

Before running the codebook over some four hundred novels we need to know what
an LLM annotation of a novel is actually worth; automated annotation without
that check is not a measurement \cite{pangakis2023validation}. We therefore
hand-annotated one novel, Lucie Delarue-Mardrus's \emph{L'Ange et les pervers}
(1930), following the same codebook we give the models. The choice is
deliberate rather than convenient: its protagonist, Marion / Mario de
Valdeclare, is presented as intersex and moves between sapphic salons and male
homosexual circles, so the novel exercises the parts of the grid we care most
about --- the \texttt{queer} value of \texttt{relation\_type}, collective and
milieu relationships, and desire that is coded rather than declared. The human
annotation records fourteen relationships, from the central bond between Marion
and Laurette Wells down to a two-line marriage recalled in the backstory. One
novel gives a first reading of how the pipeline behaves, not a corpus-level
estimate of its accuracy; everything below should be read as a pilot.

\begin{table}[h]
  \centering
  \small
  \begin{tabular}{ll}
    \toprule
    Provider & Model \\
    \midrule
    DeepSeek  & V4 Pro \\
    Anthropic & Claude Sonnet 5 \\
    Anthropic & Claude Opus 5 \\
    OpenAI    & GPT-5.6 Luna Pro \\
    OpenAI    & GPT-5.6 Terra \\
    \bottomrule
  \end{tabular}
  \caption{The five models compared, all called through OpenRouter with an
  identical system prompt, codebook, and JSON schema.}
  \label{tab:models}
\end{table}

\subsection{Two measures and their product}

Reading a novel with this grid means doing two rather different things. The
annotator first draws up a \emph{census}: which relationships in this novel are
worth an entry at all? Then, for each entry, they \emph{code} it, choosing one
label per closed field. A model can do one well and the other badly --- find
exactly the fourteen relationships a human found and misread the tone of every
one of them, or code impeccably three relationships nobody would have listed. A
single accuracy figure hides that difference, which is precisely the difference
we want to report.

The distinction is not ours. Content analysis has long separated
\emph{unitizing} --- deciding what the units of analysis are --- from
\emph{categorizing} them, and holds that the reliability of the first cannot be
folded into the reliability of the second
\cite{krippendorff2013content,krippendorff2016unitizing}. The measure that
addresses both at once is $\gamma$, which aligns annotators' units and scores
agreement on that alignment, with $\gamma_{cat}$ isolating the categorization
component \cite{mathet2015gamma,mathet2017gammacat}. We would use $\gamma$ if we
could. We cannot, for two reasons. $\gamma$ assumes units positioned on a
continuum and counts positional discrepancy as part of its dissimilarity,
whereas a relationship is not located at a place in the novel: it is identified
by \emph{who is in it}, and its footprint is the whole book. And our object is
not agreement between interchangeable observers but the distance between a
system and a reference, which is the setting precision and recall were made for.
What follows therefore keeps $\gamma$'s architecture --- align first, score
categorization on the alignment --- on a set-identity criterion rather than a
positional one.

\paragraph{Matching.} Let $G$ be the set of relationships in the human
annotation of a novel and $P$ the set proposed by a model. Neither is an ordered
list, so nothing can be compared until the two are aligned. We align on the
participants: each relationship $r$ carries an identifier $\mathrm{id}(r)$ built
from the characters involved --- \texttt{laurette\_aimee} for the affair between
Laurette Wells and Aimée de Lagres, \texttt{pere\_mere} for Marion's parents ---
recorded in both files, with aliases resolved beforehand (Marion, Mario and Miss
Marion Hervin are one character). Two relationships match when their identifiers
are equal:
%
\begin{align}
M = \{(g,p) \in G \times P \;:\; \mathrm{id}(g) = \mathrm{id}(p)\} .
\end{align}
%
In the vocabulary of coreference scoring this is CEAF under Luo's strictest
similarity $\phi_1$, which credits an alignment only when the two sets are
identical \cite{luo2005coreference}. Luo sets $\phi_1$ aside as too severe for
coreference, where exact agreement on a mention chain is vanishingly rare; it is
usable here only because participant lists are short and canonicalised upstream.
The gain is that the census score carries no similarity function and no
threshold, and a disagreement about who is involved in a relationship is treated
as what it is: a disagreement about the novel, not a spelling problem.

\paragraph{Recognition quality: did the model see the same relationships?} From
$M$ we take precision, recall and their harmonic mean, computed over
relationships rather than tokens:
%
\begin{align}
\mathrm{P} = \frac{|M|}{|P|}, \qquad
\mathrm{R} = \frac{|M|}{|G|}, \qquad
\mathrm{RQ} = \frac{2\,\mathrm{P}\,\mathrm{R}}{\mathrm{P} + \mathrm{R}} .
\end{align}
%
Precision asks whether the relationships a model proposes are ones a human
reader would also have listed; recall asks how much of the novel's relational
fabric it leaves out. RQ says nothing about how well anything was described.

\paragraph{Coding quality: did it fill the grid the same way?} For a matched
pair $(g,p)$ we compare the closed fields one by one and take the proportion on
which the human and the model wrote the same label:
%
\begin{align}
a(g,p) = \frac{1}{K} \sum_{k=1}^{K} \mathbf{1}\!\left[\, g_k = p_k \,\right] ,
\qquad
\mathrm{CQ} = \frac{1}{|M|} \sum_{(g,p) \in M} a(g,p) .
\end{align}
%
Credit is strict: the codebook offers no principled way to call
\texttt{decadent\_sublime} half-right when the annotator wrote
\texttt{moralized\_pathologized}. Every field weighs the same. The grid has
eleven closed fields, but the gold file predates \texttt{explicitness\_in\_plot},
added in schema v2, so this pilot scores $K = 10$. \texttt{confidence} is logged
by the models but is metadata about the annotation, not a property of the
relationship, and is not scored.

\paragraph{Annotation quality: one number, and why it is a product.} Suppose we
scored the whole annotation at once, with a matched relationship counting not as
$1$ but as $a(g,p)$. The harmonic mean of the resulting precision and recall is
exactly $\mathrm{RQ} \times \mathrm{CQ}$: the mean coding score factors out of
the numerator and cancels in the denominator. The product is therefore not a
weighting we chose but that single figure, read in its two parts:
%
\begin{align}
\mathrm{AQ} = \mathrm{RQ} \times \mathrm{CQ} .
\end{align}
%
A model that finds seven tenths of the relationships and gets seven tenths of
the grid right on those has produced about half the annotation we wanted, which
is what $\mathrm{AQ} \approx 0.51$ says. AQ is bounded above by RQ and reaches
it only if every matched relationship is coded perfectly. We always report the
triple $(\mathrm{RQ}, \mathrm{CQ}, \mathrm{AQ})$: two models can arrive at the
same AQ by opposite routes, one prolific and careless, the other cautious and
narrow, and for a corpus study those are not interchangeable.

This is the shape of panoptic quality in image segmentation, where a recognition
quality --- the same $F_1$ over matched objects --- multiplies a segmentation
quality averaging overlap on the matches \cite{kirillov2019panoptic}, and of
template-filling evaluation since MUC \cite{chinchor1992muc4}. We claim no new
measure: RQ, CQ and AQ instantiate a well-described family, match substructures
then normalise \cite{chen2023unified}, for a hermeneutic coding grid. What the
instantiation buys is the ability to read the two capacities apart.

\subsection{What the pipeline gets right, and where it stops}

\begin{table}[h]
  \centering
  \small
  \begin{tabular}{lcccccc}
    \toprule
    Model & $|P|$ & P & R & RQ & CQ & AQ \\
    \midrule
    \textbf{Claude Opus 5} & 14 & 0.71 & 0.71 & 0.71 & 0.72 & \textbf{0.51} \\
    GPT-5.6 Terra          &  6 & 1.00 & 0.43 & 0.60 & 0.78 & 0.47 \\
    Claude Sonnet 5        &  6 & 1.00 & 0.43 & 0.60 & 0.73 & 0.44 \\
    GPT-5.6 Luna Pro       &  5 & 1.00 & 0.36 & 0.53 & 0.66 & 0.35 \\
    DeepSeek V4 Pro        &  7 & 0.86 & 0.43 & 0.57 & 0.47 & 0.27 \\
    \bottomrule
  \end{tabular}
  \caption{Census, coding and combined score for the five models on
  \emph{L'Ange et les pervers}, against fourteen human-annotated relationships,
  sorted by AQ. $|P|$ is the number of relationships the model proposed.}
  \label{tab:model-results}
\end{table}

The $|P|$ column is the most telling. Claude Opus 5 proposes exactly fourteen
relationships and gets ten of them right; the other four propose five to seven
and are almost never wrong about the ones they propose. That second behaviour is
not caution about \emph{coding} but a narrower idea of what counts as a
relationship: those models list the named couples and stop. What they leave out
is the collective material --- Laurette Wells surrounded by her favourites, the
theatre milieu around Ginette Lobre, the powdered young men of Julien Midalge's
circle. For a study of sapphic and homosexual sociability those are not
decorative background; they are the milieux the corpus was assembled to
describe. A model with perfect precision that returns only the foreground
couples would quietly reshape the corpus into a series of couples, which is why
precision alone is a poor guide to model choice here.

\begin{table}[h]
  \centering
  \small
  \begin{tabular}{lcccccc}
    \toprule
    Field & Opus & Terra & Sonnet & Luna & DeepS. & mean \\
    \midrule
    \texttt{relation\_type} & 0.90 & 1.00 & 1.00 & 1.00 & 0.83 & 0.95 \\
    \texttt{structure}      & 0.90 & 1.00 & 0.67 & 1.00 & 1.00 & 0.91 \\
    \texttt{centrality}     & 0.90 & 0.83 & 1.00 & 0.80 & 0.67 & 0.84 \\
    \texttt{status}         & 0.70 & 1.00 & 0.83 & 0.80 & 0.83 & 0.83 \\
    \midrule
    \texttt{explicitness}   & 0.70 & 1.00 & 0.67 & 0.80 & 0.33 & 0.70 \\
    \texttt{reciprocity}    & 0.70 & 0.83 & 0.83 & 0.80 & 0.33 & 0.70 \\
    \texttt{outcome}        & 0.70 & 0.83 & 0.67 & 0.60 & 0.33 & 0.63 \\
    \midrule
    \texttt{exclusivity}    & 0.70 & 0.50 & 0.83 & 0.40 & 0.00 & 0.49 \\
    \texttt{adultery}       & 0.70 & 0.33 & 0.67 & 0.40 & 0.17 & 0.45 \\
    \texttt{framing}        & 0.30 & 0.50 & 0.17 & 0.00 & 0.17 & 0.23 \\
    \bottomrule
  \end{tabular}
  \caption{Agreement with the human annotation, field by field, on matched
  relationships only (10 pairs for Opus, 6 for Terra, Sonnet and DeepSeek, 5 for
  Luna); the last column is the unweighted mean of the five rates. Fields are
  grouped by mean agreement, not by codebook family.}
  \label{tab:field-results}
\end{table}

Averaging over fields hides the finding, so Table~\ref{tab:field-results}
reports agreement field by field, and the grid falls into three layers that are
not the codebook's three families. The top layer is descriptive and essentially
solved: once a model and a human agree that a relationship exists, they agree on
its gender configuration, the shape of the group, the kind of bond and how much
room the novel gives it. The \texttt{queer} coding is the strongest case --- every
relationship involving Marion / Mario is coded \texttt{queer} by all five models,
which is not trivial for a 1930 novel that never uses a clinical vocabulary for
its protagonist.

The middle layer sits at roughly two thirds, usable with care. The bottom layer
is where the pipeline fails, and its three fields have something in common: they
all ask for a judgment about norms. \texttt{exclusivity} and \texttt{adultery}
record the rules a relation operates under and whether they are broken, and the
codebook binds them by an explicit pairing rule; the models do not apply it as a
rule, which suggests a schema constraint rather than a better prompt.
\texttt{framing} resists most. \texttt{neutral\_ambivalent} is the models'
default, given fifteen times out of thirty-three, and
\texttt{redemptive\_spiritualized} --- which the annotator uses five times --- is
never once produced by any model. The vocabulary of \emph{pervers} and of
salvation that gives the novel its title is exactly what the automatic pass does
not recover. That contrast \emph{is} the result: the grid contains a descriptive
layer that can be automated at corpus scale today and a normative layer that
cannot, and the second carries the interpretive claims of the project.

\subsection{Limits of the measure}

Three limits bear directly on how the table should be read. \emph{Exact keying
makes collective relationships unscoreable.} Alias resolution settles naming for
named characters; it does not decide whether Laurette and her favourites are one
\texttt{hierarchy\_group} relationship or six dyads. That is a scholarly
decision, not noise, and under an exact key a legitimate disagreement about it
costs a false positive and a false negative at once. Fixing a canonical
identifier per milieu is the first change planned for the grid, alongside a
\texttt{direction} field, since the participant list stores no roles and a model
that inverts the seducer and the seduced currently scores a perfect match.

\emph{CQ is measured on a sample each model chooses for itself.} Coding quality
is computed only on matched pairs, so Terra's 0.78 averages over six legible
couples while Opus's 0.72 averages over ten, including the collective forms the
others declined --- the prolific model is penalised for attempting the hard
cases. The ordering of the middle three models in
Table~\ref{tab:model-results} should not be read as a separation.

\emph{Agreement is not chance-corrected, and there is no human ceiling.} The
marginals are skewed and the models almost never abstain where the annotator
does, so the top of the field table is flattered \cite{dieugenio2004kappa}. A
chance-corrected coefficient becomes meaningful once a second human annotator
gives us something to compare it against: fields where two readers would agree
at 0.6 cannot be used to convict a model at 0.6, and double annotation of a
subset is the measurement that would make
Table~\ref{tab:field-results} fully interpretable.

\subsection{What we carry forward}

On this pilot, Claude Opus 5 is the model we take to the rest of the corpus: it
is the only one that annotates at the granularity the codebook asks for, milieux
included, and it pays for that coverage with a precision the others never have
to risk. The ranking rests on one novel and should be revisited as the gold
standard grows.

More important than the ranking is the shape of the result. The census and the
descriptive half of the grid --- who is involved, in what gender configuration,
in what shape of group, in what kind of bond, with how much narrative room ---
are reliable enough to run over four hundred novels and to support
distributional claims. The normative half --- what the relation owes to whom,
and what the novel makes of it --- is not, and we do not intend to publish
corpus-wide figures on \texttt{framing} on the strength of an automatic pass.
The pipeline's role is a division of labour rather than a replacement: it
produces the census and the descriptive coding at scale, and directs human
reading towards the relationships where the interpretive fields are actually at
stake..
%
% Bibliography keys required:
%   pangakis2023validation, krippendorff2013content, krippendorff2016unitizing,
%   mathet2015gamma, mathet2017gammacat, luo2005coreference,
%   kirillov2019panoptic, chinchor1992muc4, chen2023unified, dieugenio2004kappa
%
% Figures are computed after validating every label in both files against the
% codebook vocabulary (one variant spelling in the gold file, "asymetrical",
% was silently costing seven points on reciprocity before this check existed).

\section{Evaluating the annotation pipeline}
\label{sec:evaluation}

\subsection{A gold standard of one novel}

Before running the codebook over some four hundred novels we need to know what
an LLM annotation of a novel is actually worth; automated annotation without
that check is not a measurement \cite{pangakis2023validation}. We therefore
hand-annotated one novel, Lucie Delarue-Mardrus's \emph{L'Ange et les pervers}
(1930), following the same codebook we give the models. The choice is
deliberate rather than convenient: its protagonist, Marion / Mario de
Valdeclare, is presented as intersex and moves between sapphic salons and male
homosexual circles, so the novel exercises the parts of the grid we care most
about --- the \texttt{queer} value of \texttt{relation\_type}, collective and
milieu relationships, and desire that is coded rather than declared. The human
annotation records fourteen relationships, from the central bond between Marion
and Laurette Wells down to a two-line marriage recalled in the backstory. One
novel gives a first reading of how the pipeline behaves, not a corpus-level
estimate of its accuracy; everything below should be read as a pilot.

\begin{table}[h]
  \centering
  \small
  \begin{tabular}{ll}
    \toprule
    Provider & Model \\
    \midrule
    DeepSeek  & V4 Pro \\
    Anthropic & Claude Sonnet 5 \\
    Anthropic & Claude Opus 5 \\
    OpenAI    & GPT-5.6 Luna Pro \\
    OpenAI    & GPT-5.6 Terra \\
    \bottomrule
  \end{tabular}
  \caption{The five models compared, all called through OpenRouter with an
  identical system prompt, codebook, and JSON schema.}
  \label{tab:models}
\end{table}

\subsection{Two measures and their product}

Reading a novel with this grid means doing two rather different things. The
annotator first draws up a \emph{census}: which relationships in this novel are
worth an entry at all? Then, for each entry, they \emph{code} it, choosing one
label per closed field. A model can do one well and the other badly --- find
exactly the fourteen relationships a human found and misread the tone of every
one of them, or code impeccably three relationships nobody would have listed. A
single accuracy figure hides that difference, which is precisely the difference
we want to report.

The distinction is not ours. Content analysis has long separated
\emph{unitizing} --- deciding what the units of analysis are --- from
\emph{categorizing} them, and holds that the reliability of the first cannot be
folded into the reliability of the second
\cite{krippendorff2013content,krippendorff2016unitizing}. The measure that
addresses both at once is $\gamma$, which aligns annotators' units and scores
agreement on that alignment, with $\gamma_{cat}$ isolating the categorization
component \cite{mathet2015gamma,mathet2017gammacat}. We would use $\gamma$ if we
could. We cannot, for two reasons. $\gamma$ assumes units positioned on a
continuum and counts positional discrepancy as part of its dissimilarity,
whereas a relationship is not located at a place in the novel: it is identified
by \emph{who is in it}, and its footprint is the whole book. And our object is
not agreement between interchangeable observers but the distance between a
system and a reference, which is the setting precision and recall were made for.
What follows therefore keeps $\gamma$'s architecture --- align first, score
categorization on the alignment --- on a set-identity criterion rather than a
positional one.

\paragraph{Matching.} Let $G$ be the set of relationships in the human
annotation of a novel and $P$ the set proposed by a model. Neither is an ordered
list, so nothing can be compared until the two are aligned. We align on the
participants: each relationship $r$ carries an identifier $\mathrm{id}(r)$ built
from the characters involved --- \texttt{laurette\_aimee} for the affair between
Laurette Wells and Aimée de Lagres, \texttt{pere\_mere} for Marion's parents ---
recorded in both files, with aliases resolved beforehand (Marion, Mario and Miss
Marion Hervin are one character). Two relationships match when their identifiers
are equal:
%
\begin{align}
M = \{(g,p) \in G \times P \;:\; \mathrm{id}(g) = \mathrm{id}(p)\} .
\end{align}
%
In the vocabulary of coreference scoring this is CEAF under Luo's strictest
similarity $\phi_1$, which credits an alignment only when the two sets are
identical \cite{luo2005coreference}. Luo sets $\phi_1$ aside as too severe for
coreference, where exact agreement on a mention chain is vanishingly rare; it is
usable here only because participant lists are short and canonicalised upstream.
The gain is that the census score carries no similarity function and no
threshold, and a disagreement about who is involved in a relationship is treated
as what it is: a disagreement about the novel, not a spelling problem.

\paragraph{Recognition quality: did the model see the same relationships?} From
$M$ we take precision, recall and their harmonic mean, computed over
relationships rather than tokens:
%
\begin{align}
\mathrm{P} = \frac{|M|}{|P|}, \qquad
\mathrm{R} = \frac{|M|}{|G|}, \qquad
\mathrm{RQ} = \frac{2\,\mathrm{P}\,\mathrm{R}}{\mathrm{P} + \mathrm{R}} .
\end{align}
%
Precision asks whether the relationships a model proposes are ones a human
reader would also have listed; recall asks how much of the novel's relational
fabric it leaves out. RQ says nothing about how well anything was described.

\paragraph{Coding quality: did it fill the grid the same way?} For a matched
pair $(g,p)$ we compare the closed fields one by one and take the proportion on
which the human and the model wrote the same label:
%
\begin{align}
a(g,p) = \frac{1}{K} \sum_{k=1}^{K} \mathbf{1}\!\left[\, g_k = p_k \,\right] ,
\qquad
\mathrm{CQ} = \frac{1}{|M|} \sum_{(g,p) \in M} a(g,p) .
\end{align}
%
Credit is strict: the codebook offers no principled way to call
\texttt{decadent\_sublime} half-right when the annotator wrote
\texttt{moralized\_pathologized}. Every field weighs the same. The grid has
eleven closed fields, but the gold file predates \texttt{explicitness\_in\_plot},
added in schema v2, so this pilot scores $K = 10$. \texttt{confidence} is logged
by the models but is metadata about the annotation, not a property of the
relationship, and is not scored.

\paragraph{Annotation quality: one number, and why it is a product.} Suppose we
scored the whole annotation at once, with a matched relationship counting not as
$1$ but as $a(g,p)$. The harmonic mean of the resulting precision and recall is
exactly $\mathrm{RQ} \times \mathrm{CQ}$: the mean coding score factors out of
the numerator and cancels in the denominator. The product is therefore not a
weighting we chose but that single figure, read in its two parts:
%
\begin{align}
\mathrm{AQ} = \mathrm{RQ} \times \mathrm{CQ} .
\end{align}
%
A model that finds seven tenths of the relationships and gets seven tenths of
the grid right on those has produced about half the annotation we wanted, which
is what $\mathrm{AQ} \approx 0.51$ says. AQ is bounded above by RQ and reaches
it only if every matched relationship is coded perfectly. We always report the
triple $(\mathrm{RQ}, \mathrm{CQ}, \mathrm{AQ})$: two models can arrive at the
same AQ by opposite routes, one prolific and careless, the other cautious and
narrow, and for a corpus study those are not interchangeable.

This is the shape of panoptic quality in image segmentation, where a recognition
quality --- the same $F_1$ over matched objects --- multiplies a segmentation
quality averaging overlap on the matches \cite{kirillov2019panoptic}, and of
template-filling evaluation since MUC \cite{chinchor1992muc4}. We claim no new
measure: RQ, CQ and AQ instantiate a well-described family, match substructures
then normalise \cite{chen2023unified}, for a hermeneutic coding grid. What the
instantiation buys is the ability to read the two capacities apart.

\subsection{What the pipeline gets right, and where it stops}

\begin{table}[h]
  \centering
  \small
  \begin{tabular}{lcccccc}
    \toprule
    Model & $|P|$ & P & R & RQ & CQ & AQ \\
    \midrule
    \textbf{Claude Opus 5} & 14 & 0.71 & 0.71 & 0.71 & 0.72 & \textbf{0.51} \\
    GPT-5.6 Terra          &  6 & 1.00 & 0.43 & 0.60 & 0.78 & 0.47 \\
    Claude Sonnet 5        &  6 & 1.00 & 0.43 & 0.60 & 0.73 & 0.44 \\
    GPT-5.6 Luna Pro       &  5 & 1.00 & 0.36 & 0.53 & 0.66 & 0.35 \\
    DeepSeek V4 Pro        &  7 & 0.86 & 0.43 & 0.57 & 0.47 & 0.27 \\
    \bottomrule
  \end{tabular}
  \caption{Census, coding and combined score for the five models on
  \emph{L'Ange et les pervers}, against fourteen human-annotated relationships,
  sorted by AQ. $|P|$ is the number of relationships the model proposed.}
  \label{tab:model-results}
\end{table}

The $|P|$ column is the most telling. Claude Opus 5 proposes exactly fourteen
relationships and gets ten of them right; the other four propose five to seven
and are almost never wrong about the ones they propose. That second behaviour is
not caution about \emph{coding} but a narrower idea of what counts as a
relationship: those models list the named couples and stop. What they leave out
is the collective material --- Laurette Wells surrounded by her favourites, the
theatre milieu around Ginette Lobre, the powdered young men of Julien Midalge's
circle. For a study of sapphic and homosexual sociability those are not
decorative background; they are the milieux the corpus was assembled to
describe. A model with perfect precision that returns only the foreground
couples would quietly reshape the corpus into a series of couples, which is why
precision alone is a poor guide to model choice here.

\begin{table}[h]
  \centering
  \small
  \begin{tabular}{lcccccc}
    \toprule
    Field & Opus & Terra & Sonnet & Luna & DeepS. & mean \\
    \midrule
    \texttt{relation\_type} & 0.90 & 1.00 & 1.00 & 1.00 & 0.83 & 0.95 \\
    \texttt{structure}      & 0.90 & 1.00 & 0.67 & 1.00 & 1.00 & 0.91 \\
    \texttt{centrality}     & 0.90 & 0.83 & 1.00 & 0.80 & 0.67 & 0.84 \\
    \texttt{status}         & 0.70 & 1.00 & 0.83 & 0.80 & 0.83 & 0.83 \\
    \midrule
    \texttt{explicitness}   & 0.70 & 1.00 & 0.67 & 0.80 & 0.33 & 0.70 \\
    \texttt{reciprocity}    & 0.70 & 0.83 & 0.83 & 0.80 & 0.33 & 0.70 \\
    \texttt{outcome}        & 0.70 & 0.83 & 0.67 & 0.60 & 0.33 & 0.63 \\
    \midrule
    \texttt{exclusivity}    & 0.70 & 0.50 & 0.83 & 0.40 & 0.00 & 0.49 \\
    \texttt{adultery}       & 0.70 & 0.33 & 0.67 & 0.40 & 0.17 & 0.45 \\
    \texttt{framing}        & 0.30 & 0.50 & 0.17 & 0.00 & 0.17 & 0.23 \\
    \bottomrule
  \end{tabular}
  \caption{Agreement with the human annotation, field by field, on matched
  relationships only (10 pairs for Opus, 6 for Terra, Sonnet and DeepSeek, 5 for
  Luna); the last column is the unweighted mean of the five rates. Fields are
  grouped by mean agreement, not by codebook family.}
  \label{tab:field-results}
\end{table}

Averaging over fields hides the finding, so Table~\ref{tab:field-results}
reports agreement field by field, and the grid falls into three layers that are
not the codebook's three families. The top layer is descriptive and essentially
solved: once a model and a human agree that a relationship exists, they agree on
its gender configuration, the shape of the group, the kind of bond and how much
room the novel gives it. The \texttt{queer} coding is the strongest case --- every
relationship involving Marion / Mario is coded \texttt{queer} by all five models,
which is not trivial for a 1930 novel that never uses a clinical vocabulary for
its protagonist.

The middle layer sits at roughly two thirds, usable with care. The bottom layer
is where the pipeline fails, and its three fields have something in common: they
all ask for a judgment about norms. \texttt{exclusivity} and \texttt{adultery}
record the rules a relation operates under and whether they are broken, and the
codebook binds them by an explicit pairing rule; the models do not apply it as a
rule, which suggests a schema constraint rather than a better prompt.
\texttt{framing} resists most. \texttt{neutral\_ambivalent} is the models'
default, given fifteen times out of thirty-three, and
\texttt{redemptive\_spiritualized} --- which the annotator uses five times --- is
never once produced by any model. The vocabulary of \emph{pervers} and of
salvation that gives the novel its title is exactly what the automatic pass does
not recover. That contrast \emph{is} the result: the grid contains a descriptive
layer that can be automated at corpus scale today and a normative layer that
cannot, and the second carries the interpretive claims of the project.

\subsection{Limits of the measure}

Three limits bear directly on how the table should be read. \emph{Exact keying
makes collective relationships unscoreable.} Alias resolution settles naming for
named characters; it does not decide whether Laurette and her favourites are one
\texttt{hierarchy\_group} relationship or six dyads. That is a scholarly
decision, not noise, and under an exact key a legitimate disagreement about it
costs a false positive and a false negative at once. Fixing a canonical
identifier per milieu is the first change planned for the grid, alongside a
\texttt{direction} field, since the participant list stores no roles and a model
that inverts the seducer and the seduced currently scores a perfect match.

\emph{CQ is measured on a sample each model chooses for itself.} Coding quality
is computed only on matched pairs, so Terra's 0.78 averages over six legible
couples while Opus's 0.72 averages over ten, including the collective forms the
others declined --- the prolific model is penalised for attempting the hard
cases. The ordering of the middle three models in
Table~\ref{tab:model-results} should not be read as a separation.

\emph{Agreement is not chance-corrected, and there is no human ceiling.} The
marginals are skewed and the models almost never abstain where the annotator
does, so the top of the field table is flattered \cite{dieugenio2004kappa}. A
chance-corrected coefficient becomes meaningful once a second human annotator
gives us something to compare it against: fields where two readers would agree
at 0.6 cannot be used to convict a model at 0.6, and double annotation of a
subset is the measurement that would make
Table~\ref{tab:field-results} fully interpretable.

\subsection{What we carry forward}

On this pilot, Claude Opus 5 is the model we take to the rest of the corpus: it
is the only one that annotates at the granularity the codebook asks for, milieux
included, and it pays for that coverage with a precision the others never have
to risk. The ranking rests on one novel and should be revisited as the gold
standard grows.

More important than the ranking is the shape of the result. The census and the
descriptive half of the grid --- who is involved, in what gender configuration,
in what shape of group, in what kind of bond, with how much narrative room ---
are reliable enough to run over four hundred novels and to support
distributional claims. The normative half --- what the relation owes to whom,
and what the novel makes of it --- is not, and we do not intend to publish
corpus-wide figures on \texttt{framing} on the strength of an automatic pass.
The pipeline's role is a division of labour rather than a replacement: it
produces the census and the descriptive coding at scale, and directs human
reading towards the relationships where the interpretive fields are actually at
stake..
%
% Bibliography keys required:
%   pangakis2023validation, krippendorff2013content, krippendorff2016unitizing,
%   mathet2015gamma, mathet2017gammacat, luo2005coreference,
%   kirillov2019panoptic, chinchor1992muc4, chen2023unified, dieugenio2004kappa
%
% Figures are computed after validating every label in both files against the
% codebook vocabulary (one variant spelling in the gold file, "asymetrical",
% was silently costing seven points on reciprocity before this check existed).

\section{Evaluating the annotation pipeline}
\label{sec:evaluation}

\subsection{A gold standard of one novel}

Before running the codebook over some four hundred novels we need to know what
an LLM annotation of a novel is actually worth; automated annotation without
that check is not a measurement \cite{pangakis2023validation}. We therefore
hand-annotated one novel, Lucie Delarue-Mardrus's \emph{L'Ange et les pervers}
(1930), following the same codebook we give the models. The choice is
deliberate rather than convenient: its protagonist, Marion / Mario de
Valdeclare, is presented as intersex and moves between sapphic salons and male
homosexual circles, so the novel exercises the parts of the grid we care most
about --- the \texttt{queer} value of \texttt{relation\_type}, collective and
milieu relationships, and desire that is coded rather than declared. The human
annotation records fourteen relationships, from the central bond between Marion
and Laurette Wells down to a two-line marriage recalled in the backstory. One
novel gives a first reading of how the pipeline behaves, not a corpus-level
estimate of its accuracy; everything below should be read as a pilot.

\begin{table}[h]
  \centering
  \small
  \begin{tabular}{ll}
    \toprule
    Provider & Model \\
    \midrule
    DeepSeek  & V4 Pro \\
    Anthropic & Claude Sonnet 5 \\
    Anthropic & Claude Opus 5 \\
    OpenAI    & GPT-5.6 Luna Pro \\
    OpenAI    & GPT-5.6 Terra \\
    \bottomrule
  \end{tabular}
  \caption{The five models compared, all called through OpenRouter with an
  identical system prompt, codebook, and JSON schema.}
  \label{tab:models}
\end{table}

\subsection{Two measures and their product}

Reading a novel with this grid means doing two rather different things. The
annotator first draws up a \emph{census}: which relationships in this novel are
worth an entry at all? Then, for each entry, they \emph{code} it, choosing one
label per closed field. A model can do one well and the other badly --- find
exactly the fourteen relationships a human found and misread the tone of every
one of them, or code impeccably three relationships nobody would have listed. A
single accuracy figure hides that difference, which is precisely the difference
we want to report.

The distinction is not ours. Content analysis has long separated
\emph{unitizing} --- deciding what the units of analysis are --- from
\emph{categorizing} them, and holds that the reliability of the first cannot be
folded into the reliability of the second
\cite{krippendorff2013content,krippendorff2016unitizing}. The measure that
addresses both at once is $\gamma$, which aligns annotators' units and scores
agreement on that alignment, with $\gamma_{cat}$ isolating the categorization
component \cite{mathet2015gamma,mathet2017gammacat}. We would use $\gamma$ if we
could. We cannot, for two reasons. $\gamma$ assumes units positioned on a
continuum and counts positional discrepancy as part of its dissimilarity,
whereas a relationship is not located at a place in the novel: it is identified
by \emph{who is in it}, and its footprint is the whole book. And our object is
not agreement between interchangeable observers but the distance between a
system and a reference, which is the setting precision and recall were made for.
What follows therefore keeps $\gamma$'s architecture --- align first, score
categorization on the alignment --- on a set-identity criterion rather than a
positional one.

\paragraph{Matching.} Let $G$ be the set of relationships in the human
annotation of a novel and $P$ the set proposed by a model. Neither is an ordered
list, so nothing can be compared until the two are aligned. We align on the
participants: each relationship $r$ carries an identifier $\mathrm{id}(r)$ built
from the characters involved --- \texttt{laurette\_aimee} for the affair between
Laurette Wells and Aimée de Lagres, \texttt{pere\_mere} for Marion's parents ---
recorded in both files, with aliases resolved beforehand (Marion, Mario and Miss
Marion Hervin are one character). Two relationships match when their identifiers
are equal:
%
\begin{align}
M = \{(g,p) \in G \times P \;:\; \mathrm{id}(g) = \mathrm{id}(p)\} .
\end{align}
%
In the vocabulary of coreference scoring this is CEAF under Luo's strictest
similarity $\phi_1$, which credits an alignment only when the two sets are
identical \cite{luo2005coreference}. Luo sets $\phi_1$ aside as too severe for
coreference, where exact agreement on a mention chain is vanishingly rare; it is
usable here only because participant lists are short and canonicalised upstream.
The gain is that the census score carries no similarity function and no
threshold, and a disagreement about who is involved in a relationship is treated
as what it is: a disagreement about the novel, not a spelling problem.

\paragraph{Recognition quality: did the model see the same relationships?} From
$M$ we take precision, recall and their harmonic mean, computed over
relationships rather than tokens:
%
\begin{align}
\mathrm{P} = \frac{|M|}{|P|}, \qquad
\mathrm{R} = \frac{|M|}{|G|}, \qquad
\mathrm{RQ} = \frac{2\,\mathrm{P}\,\mathrm{R}}{\mathrm{P} + \mathrm{R}} .
\end{align}
%
Precision asks whether the relationships a model proposes are ones a human
reader would also have listed; recall asks how much of the novel's relational
fabric it leaves out. RQ says nothing about how well anything was described.

\paragraph{Coding quality: did it fill the grid the same way?} For a matched
pair $(g,p)$ we compare the closed fields one by one and take the proportion on
which the human and the model wrote the same label:
%
\begin{align}
a(g,p) = \frac{1}{K} \sum_{k=1}^{K} \mathbf{1}\!\left[\, g_k = p_k \,\right] ,
\qquad
\mathrm{CQ} = \frac{1}{|M|} \sum_{(g,p) \in M} a(g,p) .
\end{align}
%
Credit is strict: the codebook offers no principled way to call
\texttt{decadent\_sublime} half-right when the annotator wrote
\texttt{moralized\_pathologized}. Every field weighs the same. The grid has
eleven closed fields, but the gold file predates \texttt{explicitness\_in\_plot},
added in schema v2, so this pilot scores $K = 10$. \texttt{confidence} is logged
by the models but is metadata about the annotation, not a property of the
relationship, and is not scored.

\paragraph{Annotation quality: one number, and why it is a product.} Suppose we
scored the whole annotation at once, with a matched relationship counting not as
$1$ but as $a(g,p)$. The harmonic mean of the resulting precision and recall is
exactly $\mathrm{RQ} \times \mathrm{CQ}$: the mean coding score factors out of
the numerator and cancels in the denominator. The product is therefore not a
weighting we chose but that single figure, read in its two parts:
%
\begin{align}
\mathrm{AQ} = \mathrm{RQ} \times \mathrm{CQ} .
\end{align}
%
A model that finds seven tenths of the relationships and gets seven tenths of
the grid right on those has produced about half the annotation we wanted, which
is what $\mathrm{AQ} \approx 0.51$ says. AQ is bounded above by RQ and reaches
it only if every matched relationship is coded perfectly. We always report the
triple $(\mathrm{RQ}, \mathrm{CQ}, \mathrm{AQ})$: two models can arrive at the
same AQ by opposite routes, one prolific and careless, the other cautious and
narrow, and for a corpus study those are not interchangeable.

This is the shape of panoptic quality in image segmentation, where a recognition
quality --- the same $F_1$ over matched objects --- multiplies a segmentation
quality averaging overlap on the matches \cite{kirillov2019panoptic}, and of
template-filling evaluation since MUC \cite{chinchor1992muc4}. We claim no new
measure: RQ, CQ and AQ instantiate a well-described family, match substructures
then normalise \cite{chen2023unified}, for a hermeneutic coding grid. What the
instantiation buys is the ability to read the two capacities apart.

\subsection{What the pipeline gets right, and where it stops}

\begin{table}[h]
  \centering
  \small
  \begin{tabular}{lcccccc}
    \toprule
    Model & $|P|$ & P & R & RQ & CQ & AQ \\
    \midrule
    \textbf{Claude Opus 5} & 14 & 0.71 & 0.71 & 0.71 & 0.72 & \textbf{0.51} \\
    GPT-5.6 Terra          &  6 & 1.00 & 0.43 & 0.60 & 0.78 & 0.47 \\
    Claude Sonnet 5        &  6 & 1.00 & 0.43 & 0.60 & 0.73 & 0.44 \\
    GPT-5.6 Luna Pro       &  5 & 1.00 & 0.36 & 0.53 & 0.66 & 0.35 \\
    DeepSeek V4 Pro        &  7 & 0.86 & 0.43 & 0.57 & 0.47 & 0.27 \\
    \bottomrule
  \end{tabular}
  \caption{Census, coding and combined score for the five models on
  \emph{L'Ange et les pervers}, against fourteen human-annotated relationships,
  sorted by AQ. $|P|$ is the number of relationships the model proposed.}
  \label{tab:model-results}
\end{table}

The $|P|$ column is the most telling. Claude Opus 5 proposes exactly fourteen
relationships and gets ten of them right; the other four propose five to seven
and are almost never wrong about the ones they propose. That second behaviour is
not caution about \emph{coding} but a narrower idea of what counts as a
relationship: those models list the named couples and stop. What they leave out
is the collective material --- Laurette Wells surrounded by her favourites, the
theatre milieu around Ginette Lobre, the powdered young men of Julien Midalge's
circle. For a study of sapphic and homosexual sociability those are not
decorative background; they are the milieux the corpus was assembled to
describe. A model with perfect precision that returns only the foreground
couples would quietly reshape the corpus into a series of couples, which is why
precision alone is a poor guide to model choice here.

\begin{table}[h]
  \centering
  \small
  \begin{tabular}{lcccccc}
    \toprule
    Field & Opus & Terra & Sonnet & Luna & DeepS. & mean \\
    \midrule
    \texttt{relation\_type} & 0.90 & 1.00 & 1.00 & 1.00 & 0.83 & 0.95 \\
    \texttt{structure}      & 0.90 & 1.00 & 0.67 & 1.00 & 1.00 & 0.91 \\
    \texttt{centrality}     & 0.90 & 0.83 & 1.00 & 0.80 & 0.67 & 0.84 \\
    \texttt{status}         & 0.70 & 1.00 & 0.83 & 0.80 & 0.83 & 0.83 \\
    \midrule
    \texttt{explicitness}   & 0.70 & 1.00 & 0.67 & 0.80 & 0.33 & 0.70 \\
    \texttt{reciprocity}    & 0.70 & 0.83 & 0.83 & 0.80 & 0.33 & 0.70 \\
    \texttt{outcome}        & 0.70 & 0.83 & 0.67 & 0.60 & 0.33 & 0.63 \\
    \midrule
    \texttt{exclusivity}    & 0.70 & 0.50 & 0.83 & 0.40 & 0.00 & 0.49 \\
    \texttt{adultery}       & 0.70 & 0.33 & 0.67 & 0.40 & 0.17 & 0.45 \\
    \texttt{framing}        & 0.30 & 0.50 & 0.17 & 0.00 & 0.17 & 0.23 \\
    \bottomrule
  \end{tabular}
  \caption{Agreement with the human annotation, field by field, on matched
  relationships only (10 pairs for Opus, 6 for Terra, Sonnet and DeepSeek, 5 for
  Luna); the last column is the unweighted mean of the five rates. Fields are
  grouped by mean agreement, not by codebook family.}
  \label{tab:field-results}
\end{table}

Averaging over fields hides the finding, so Table~\ref{tab:field-results}
reports agreement field by field, and the grid falls into three layers that are
not the codebook's three families. The top layer is descriptive and essentially
solved: once a model and a human agree that a relationship exists, they agree on
its gender configuration, the shape of the group, the kind of bond and how much
room the novel gives it. The \texttt{queer} coding is the strongest case --- every
relationship involving Marion / Mario is coded \texttt{queer} by all five models,
which is not trivial for a 1930 novel that never uses a clinical vocabulary for
its protagonist.

The middle layer sits at roughly two thirds, usable with care. The bottom layer
is where the pipeline fails, and its three fields have something in common: they
all ask for a judgment about norms. \texttt{exclusivity} and \texttt{adultery}
record the rules a relation operates under and whether they are broken, and the
codebook binds them by an explicit pairing rule; the models do not apply it as a
rule, which suggests a schema constraint rather than a better prompt.
\texttt{framing} resists most. \texttt{neutral\_ambivalent} is the models'
default, given fifteen times out of thirty-three, and
\texttt{redemptive\_spiritualized} --- which the annotator uses five times --- is
never once produced by any model. The vocabulary of \emph{pervers} and of
salvation that gives the novel its title is exactly what the automatic pass does
not recover. That contrast \emph{is} the result: the grid contains a descriptive
layer that can be automated at corpus scale today and a normative layer that
cannot, and the second carries the interpretive claims of the project.

\subsection{Limits of the measure}

Three limits bear directly on how the table should be read. \emph{Exact keying
makes collective relationships unscoreable.} Alias resolution settles naming for
named characters; it does not decide whether Laurette and her favourites are one
\texttt{hierarchy\_group} relationship or six dyads. That is a scholarly
decision, not noise, and under an exact key a legitimate disagreement about it
costs a false positive and a false negative at once. Fixing a canonical
identifier per milieu is the first change planned for the grid, alongside a
\texttt{direction} field, since the participant list stores no roles and a model
that inverts the seducer and the seduced currently scores a perfect match.

\emph{CQ is measured on a sample each model chooses for itself.} Coding quality
is computed only on matched pairs, so Terra's 0.78 averages over six legible
couples while Opus's 0.72 averages over ten, including the collective forms the
others declined --- the prolific model is penalised for attempting the hard
cases. The ordering of the middle three models in
Table~\ref{tab:model-results} should not be read as a separation.

\emph{Agreement is not chance-corrected, and there is no human ceiling.} The
marginals are skewed and the models almost never abstain where the annotator
does, so the top of the field table is flattered \cite{dieugenio2004kappa}. A
chance-corrected coefficient becomes meaningful once a second human annotator
gives us something to compare it against: fields where two readers would agree
at 0.6 cannot be used to convict a model at 0.6, and double annotation of a
subset is the measurement that would make
Table~\ref{tab:field-results} fully interpretable.

\subsection{What we carry forward}

On this pilot, Claude Opus 5 is the model we take to the rest of the corpus: it
is the only one that annotates at the granularity the codebook asks for, milieux
included, and it pays for that coverage with a precision the others never have
to risk. The ranking rests on one novel and should be revisited as the gold
standard grows.

More important than the ranking is the shape of the result. The census and the
descriptive half of the grid --- who is involved, in what gender configuration,
in what shape of group, in what kind of bond, with how much narrative room ---
are reliable enough to run over four hundred novels and to support
distributional claims. The normative half --- what the relation owes to whom,
and what the novel makes of it --- is not, and we do not intend to publish
corpus-wide figures on \texttt{framing} on the strength of an automatic pass.
The pipeline's role is a division of labour rather than a replacement: it
produces the census and the descriptive coding at scale, and directs human
reading towards the relationships where the interpretive fields are actually at
stake.